
\section{Discussion}\label{sec:discussion}



\subsection{Design principles and simulation paradigms}

Simulations are proliferating in high-stakes engineering disciplines, where iterative design and validation are performed long before prototypes are built. 
In contrast, traditional mechanism design research often relies on expensive, time-intensive human subject experiments or confounded observational data. 
Our synthetic lab leverages the rapidly advancing capabilities of large language models to simulate bidder behavior at a fraction of the cost. 
By systematically varying auction parameters, rules, and descriptive framings in a controlled setting, we not only replicated well-known empirical regularities in sealed-bid second-price auctions and eBay-type proxy auctions, crucially, reveal how seemingly minor descriptive choices—while invisible in standard dominant-strategy theory—produce large shifts in over- and under-bidding behavior in practice.


In doing so, this
simulation paradigm complements traditional analytical models by integrating behavioral factors into mechanism design.

LLM agents operate on a text-based reasoning framework that captures nuanced decision-making processes. 
This also affords  the opportunity to query these  LLM agents about their internal deliberations, opening up the prospect, with further study and validation, that LLMs may also provide information on the human cognitive constraints that  underlie deviations from classical auction theory. 


\subsection{The advantage of LLM participants}
\paragraph{Relative to Traditional Lab Experiments}
Classical lab experiments in auction settings have yielded valuable insights, 
%for instance, studies of \citet{li2017obviously} employing over 400 human subjects have demonstrated 
but also come with high cost and logistical challenges.
In stark contrast, our experiments with LLM agents reduce these costs by roughly three orders of magnitude (e.g., from over $\$15{,}000$ to  $\$10$ per experiment \citep{li2017obviously}). 
% \dcp{need cite for 15000}\kz{added citation}
%
This dramatic improvement in cost efficiency enables rapid iteration, allowing researchers to explore a broader parameter space with minimal investment of money and time. 

\paragraph{Relative to Observational Data}
The analysis of observational datasets from platforms like eBay is often challenging because one
needs to control for the heterogeneity among sellers and buyers, this complicating causal inference.
%
In our synthetic lab, we can isolate individual auction parameters and rule framings by carefully 
controlling the environment in which bids are generated including the simulated goods and valuations of bidders. 
This controlled setting is perfect for
causal inference and enables us to directly link subtle changes in auction design to variations in bidder behavior \citep{manning2024automated}.


\paragraph{Online Laboratory Framework}
The rise of online human labs \citep{horton2011online} in behavioral economics has previously lowered the barrier to performing controlled experiments \citep{oprea2024decisions, enke2024behavioral}, enabling near real-time data collection and hypothesis testing. 
Our synthetic lab extends this idea by eliminating many of the ethical, logistical, and opportunity-cost issues associated with human experiments. 
By using LLM agents, we gain the flexibility to incorporate additional modules (such as reasoning queries) into the experimental design without risking harm to human subjects. 
This approach accelerates research and also offers a platform for cross-validating findings with more traditional online lab experiments.

% \subsection{Comparison with Traditional Auction Test Suites}

% Existing tools such as the {\em Combinatorial Auction Test Suite (CATS)}~\citep{leyton2000towards}  focus on  providing valuation distributions on which to evaluate the efficiency and computational performance of winner-determination and pricing algorithms.  While these suites provide a benchmark
% set of instances with which to study algorithmic performance and auction
% design  they
% do not seek to model bidder behavior. In this sense, the use of LLMs as proxy models for human bidders
% complements these test suites, offering  a means to probe the behavioral dimensions 
% bidder behaviors in auction design.
